% uid: ""
% school: "lycee"
% level: "3eme"
% section: ["Economie et Gestion"]
% chapter: "Suites réelles"
% title: "Écritures décimales périodiques et fractions"
% difficulty: "Moyen"
% points: null
% tags: ["suites", "somme-geometrique", "limite", "decimal-periodique", "fraction"]
% provenance: "own"
% dressing: "somme infinie"
% rung: 3
% batch: "2026-11-20-suites-3eco-A"

% ---STATEMENT---
\begin{enumerate}
  \item On pose $S_n=0{,}18+0{,}0018+0{,}000018+\dots$ où le $n$-ième terme (pour $n\ge0$) est $0{,}18\times(0{,}01)^n$.
  \begin{enumerate}
    \item Calculer $S_n=\sum_{k=0}^{n}0{,}18\times(0{,}01)^k$ en fonction de $n$.
    \item En déduire la limite de $S_n$ et montrer que $0{,}181818\ldots=\dfrac2{11}$.
  \end{enumerate}
  \item Écrire de même $0{,}2777\ldots$ sous forme de fraction irréductible. On pourra écrire $0{,}2777\ldots=0{,}2+0{,}0777\ldots$
\end{enumerate}

% ---SOLUTION---
\begin{enumerate}
  \item
  \begin{enumerate}
    \item \textbf{Rappel.}\par
  \begin{prop}[Somme géométrique]
    Si $q\neq1$ : $u_0+u_1+\dots+u_n=u_0\,\dfrac{1-q^{\,n+1}}{1-q}$ (il y a $n+1$ termes).
    \end{prop}
    Ici $u_0=0{,}18$ et $q=0{,}01$ :
    $$S_n=0{,}18\times\frac{1-0{,}01^{\,n+1}}{1-0{,}01}=\frac{0{,}18}{0{,}99}\bigl(1-0{,}01^{\,n+1}\bigr)=\ib{$\dfrac2{11}\bigl(1-0{,}01^{\,n+1}\bigr)$}.$$
    (On a simplifié $\dfrac{0{,}18}{0{,}99}=\dfrac{18}{99}=\dfrac2{11}$.)
    \item Comme $-1<0{,}01<1$, $0{,}01^{\,n+1}\to0$, donc $S_n\to\dfrac2{11}$.

    Le développement décimal $0{,}1818\ldots$ est exactement cette somme « infinie » : $0{,}181818\ldots=$ \ib{$\dfrac2{11}$}.

    \emph{Vérification} : $2\div11=0{,}181818\ldots$ $\checkmark$
  \end{enumerate}
  \item $0{,}2777\ldots=0{,}2+0{,}07+0{,}007+0{,}0007+\dots$

  Les termes à partir de $0{,}07$ forment une somme géométrique de premier terme $0{,}07$ et de raison $0{,}1$ :
  \begin{prop}[Somme infinie]
  Si $-1<q<1$, $u_0+u_0q+\dots+u_0q^n\to\dfrac{u_0}{1-q}$.
  \end{prop}
  $$0{,}07+0{,}007+\dots=\frac{0{,}07}{1-0{,}1}=\frac{0{,}07}{0{,}9}=\frac7{90}.$$
  $$0{,}2777\ldots=\frac15+\frac7{90}=\frac{18}{90}+\frac7{90}=\frac{25}{90}=\ib{$\dfrac5{18}$}.$$
  \emph{Vérification} : $5\div18=0{,}27777\ldots$ $\checkmark$

  \textbf{Piège :} la raison est $0{,}1$ (et non $0{,}01$) pour $0{,}0777\ldots$ : chaque terme est dix fois plus petit que le précédent.
\end{enumerate}

\textbf{Idée :} une écriture décimale périodique est la somme infinie d'une suite géométrique de raison comprise entre $0$ et $1$.
